\begin{resumen}
	La implementación del sistema web para la gestión de ESAVI se fundamenta en una arquitectura modular basada en los principios de Clean Architecture, utilizando .NET 8 para el backend y React con TypeScript para el frontend. 
	La solución integra mecanismos avanzados de seguridad, incluyendo autenticación mediante JWT con rotación de refresh tokens almacenados en cookies HttpOnly, limitación de velocidad (rate limiting) y protección contra ataques automatizados mediante Friendly CAPTCHA. 
	Para optimizar la experiencia del usuario y la fiabilidad, se implementó una comunicación asíncrona empleando RabbitMQ y MassTransit, facilitando notificaciones automáticas vía correo electrónico y WhatsApp. 
	El sistema garantiza la integridad y trazabilidad de los datos de farmacovigilancia mediante el uso del patrón Unit of Work y un registro de auditoría automatizado sobre una base de datos MySQL.
\end{resumen}

\begin{abstract}
	The implementation of the ESAVI management web system is based on a modular architecture following Clean Architecture principles. 
	The backend was developed using .NET 8 and C\#, while the frontend was built as a Single Page Application (SPA) with React and TypeScript. 
	The system features a robust security framework, including JWT authentication with Refresh Token rotation (HttpOnly cookies), IP-based rate limiting, and Friendly CAPTCHA protection for public endpoints. 
	To ensure high availability and responsiveness, the solution incorporates asynchronous messaging via RabbitMQ and MassTransit for automated email and WhatsApp notifications. 
	Data integrity and traceability are managed through Entity Framework Core, the Unit of Work pattern, and an automated Audit Log system.
\end{abstract}